% GENERATED FILE. Edit canonical lessons, then run npm run generate:books.
% canonical-source-hash: fnv1a64:d092e34f2c6a3077
% canonical-lessons: ES-C760-disfraz, ES-C760-disfrazarse, ES-C760-maquillaje, ES-C760-estreno, ES-C760-interpretar

\chapter{Teatro y disfraces}
\label{ch:es-teatro-y-disfraces}
\hlchaptermodality{760}

\begin{chapteropening}
\textbf{By the end of this chapter:} \emph{I can talk about amateur theatre and dressing up: costumes, make-up and opening night.}

Complete the last lesson of chapter 760: I can talk about amateur theatre and dressing up: costumes, make-up and opening night.
\end{chapteropening}

\section[el disfraz]{el disfraz — the costume, the disguise}

\label{lesson:ES-C760-disfraz}



\begin{quote}
Before the new one: say the Spanish for to act, then the Spanish for the performance.
\end{quote}

\subsection*{You'll want to know: el disfraz}
\textbf{el disfraz} — \textquotedblleft{}the costume, the disguise\textquotedblright{}.

\begin{quote}
Mi disfraz de carnaval es de pirata.
\end{quote}

\begin{quote}
\emph{My carnival costume is a pirate costume.}
\end{quote}

\emph{Una fiesta de disfraces} — a fancy-dress party.

\begin{cousinweb}
\emph{disfraz} is built on \emph{disfrazar} (EtymDB), from Old Spanish \emph{desfrezar}. For \emph{desfrezar} the origin is not settled.
\end{cousinweb}

\begin{grammarlens}[title={What you've built}]
One more word for telling, explaining and describing at B1.
\end{grammarlens}

\subsection*{Your turn}
\begin{itemize}
\raggedright
  \item \emph{Say it:} \emph{el disfraz}
  \item \emph{Say it:} \emph{el disfraz}, once more
  \item \emph{Say it:} a sentence of your own with \emph{el disfraz}
  \item \emph{Recall:} say the Spanish for to act, then the Spanish for the performance, then say \emph{el disfraz} again
\end{itemize}

\begin{tcolorbox}[breakable,colback=teal!4,colframe=teal!35!black,title={Before you move on}]
What does el disfraz mean? (\textquotedblleft{}The costume, the disguise\textquotedblright{}.) Say it once more.
\end{tcolorbox}
\section[disfrazarse]{disfrazarse — to dress up, to disguise oneself}

\label{lesson:ES-C760-disfrazarse}



\begin{quote}
Before the new one: say the Spanish for the performance, then the Spanish for the costume.
\end{quote}

\subsection*{You'll want to know: disfrazarse}
\textbf{disfrazarse} — \textquotedblleft{}to dress up, to disguise oneself\textquotedblright{}.

\begin{quote}
En carnaval nos disfrazamos.
\end{quote}

\begin{quote}
\emph{At carnival we dress up.}
\end{quote}

\emph{Disfrazarse de} — to dress up as.

\begin{cousinweb}
\emph{disfrazarse} comes from Old Spanish \emph{desfrezar} \textquotedblleft{}to disguise\textquotedblright{}. For \emph{desfrezar} itself the origin is not settled.
\end{cousinweb}

\begin{grammarlens}[title={What you've built}]
One more word for telling, explaining and describing at B1.
\end{grammarlens}

\subsection*{Your turn}
\begin{itemize}
\raggedright
  \item \emph{Say it:} \emph{disfrazarse}
  \item \emph{Say it:} \emph{disfrazarse}, once more
  \item \emph{Say it:} a sentence of your own with \emph{disfrazarse}
  \item \emph{Recall:} say the Spanish for the performance, then the Spanish for the costume, then say \emph{disfrazarse} again
\end{itemize}

\begin{tcolorbox}[breakable,colback=teal!4,colframe=teal!35!black,title={Before you move on}]
What does disfrazarse mean? (\textquotedblleft{}To dress up, to disguise oneself\textquotedblright{}.) Say it once more.
\end{tcolorbox}
\section[el maquillaje]{el maquillaje — make-up}

\label{lesson:ES-C760-maquillaje}



\begin{quote}
Before the new one: say the Spanish for the costume, then the Spanish for to dress up.
\end{quote}

\subsection*{You'll want to know: el maquillaje}
\textbf{el maquillaje} — \textquotedblleft{}make-up\textquotedblright{}.

\begin{quote}
Los actores necesitan mucho maquillaje.
\end{quote}

\begin{quote}
\emph{The actors need a lot of make-up.}
\end{quote}

\emph{Maquillarse} — to put on make-up.

\begin{cousinweb}
\emph{maquillaje} was borrowed from French \emph{maquillage}, built on \emph{maquiller} \textquotedblleft{}to make up\textquotedblright{}, which is usually traced to a northern French dialect verb meaning to work, to do.
\end{cousinweb}

\begin{grammarlens}[title={What you've built}]
One more word for telling, explaining and describing at B1.
\end{grammarlens}

\subsection*{Your turn}
\begin{itemize}
\raggedright
  \item \emph{Say it:} \emph{el maquillaje}
  \item \emph{Say it:} \emph{el maquillaje}, once more
  \item \emph{Say it:} a sentence of your own with \emph{el maquillaje}
  \item \emph{Recall:} say the Spanish for the costume, then the Spanish for to dress up, then say \emph{el maquillaje} again
\end{itemize}

\begin{tcolorbox}[breakable,colback=teal!4,colframe=teal!35!black,title={Before you move on}]
What does el maquillaje mean? (\textquotedblleft{}Make-up\textquotedblright{}.) Say it once more.
\end{tcolorbox}
\section[el estreno]{el estreno — the premiere, the first showing}

\label{lesson:ES-C760-estreno}



\begin{quote}
Before the new one: say the Spanish for to dress up, then the Spanish for make-up.
\end{quote}

\subsection*{You'll want to know: el estreno}
\textbf{el estreno} — \textquotedblleft{}the premiere, the first showing\textquotedblright{}.

\begin{quote}
Fuimos al estreno de la película.
\end{quote}

\begin{quote}
\emph{We went to the premiere of the film.}
\end{quote}

\emph{Estrenar} — to show for the first time, to wear for the first time.

\begin{cousinweb}
\emph{estreno} is built on \emph{estrenar}, which is usually traced to Latin \emph{strēna} \textquotedblleft{}an omen; a New Year gift\textquotedblright{} (Lewis \& Short): using something new for the first time.
\end{cousinweb}

\begin{grammarlens}[title={What you've built}]
One more word for telling, explaining and describing at B1.
\end{grammarlens}

\subsection*{Your turn}
\begin{itemize}
\raggedright
  \item \emph{Say it:} \emph{el estreno}
  \item \emph{Say it:} \emph{el estreno}, once more
  \item \emph{Say it:} a sentence of your own with \emph{el estreno}
  \item \emph{Recall:} say the Spanish for to dress up, then the Spanish for make-up, then say \emph{el estreno} again
\end{itemize}

\begin{tcolorbox}[breakable,colback=teal!4,colframe=teal!35!black,title={Before you move on}]
What does el estreno mean? (\textquotedblleft{}The premiere, the first showing\textquotedblright{}.) Say it once more.
\end{tcolorbox}
\section[interpretar]{interpretar — to play (a role); to interpret}

\label{lesson:ES-C760-interpretar}



\begin{quote}
Before the new one: say the Spanish for make-up, then the Spanish for the premiere.
\end{quote}

\subsection*{You'll want to know: interpretar}
\textbf{interpretar} — \textquotedblleft{}to play (a role); to interpret\textquotedblright{}.

\begin{quote}
Interpreta a un médico en la película.
\end{quote}

\begin{quote}
\emph{He plays a doctor in the film.}
\end{quote}

Also to interpret a language or a text.

\begin{cousinweb}
\emph{interpretar} is usually described as a learned borrowing from Latin \emph{interpretārī} \textquotedblleft{}to explain\textquotedblright{}, built on \emph{interpres} \textquotedblleft{}go-between, translator\textquotedblright{}.
\end{cousinweb}

\begin{grammarlens}[title={What you've built}]
One more word for telling, explaining and describing at B1.
\end{grammarlens}

\subsection*{Your turn}
\begin{itemize}
\raggedright
  \item \emph{Say it:} \emph{interpretar}
  \item \emph{Say it:} \emph{interpretar}, once more
  \item \emph{Say it:} a sentence of your own with \emph{interpretar}
  \item \emph{Recall:} say the Spanish for make-up, then the Spanish for the premiere, then say \emph{interpretar} again
\end{itemize}

\begin{tcolorbox}[breakable,colback=teal!4,colframe=teal!35!black,title={Before you move on}]
What does interpretar mean? (\textquotedblleft{}To play (a role); to interpret\textquotedblright{}.) Say it once more.
\end{tcolorbox}
